\cleardoublepage % memastikan bab baru dimulai di halaman ganjil (kanan)
\chapter{PENDAHULUAN}
\label{chap:pendahuluan}
% --- Latar Belakang ---
\section{Latar Belakang}
\label{sec:latar_belakang}
PT Nutech Integrasi merupakan perusahaan yang bergerak dalam penyediaan solusi teknologi informasi dan integrasi sistem untuk mendukung transformasi digital di berbagai sektor. Perusahaan memiliki empat pilar bisnis utama yang dikelompokkan ke dalam dua fokus layanan, yaitu \textit{Intelligent Mobility} (\textit{Intelligent Transportation System} dan \textit{Intelligent Logistic System}) serta \textit{Intelligent Services} (\textit{Intelligent Information System} dan \textit{Geographical Information System}). Salah satu sektor utama perusahaan adalah perkeretaapian, dengan fokus pada pengembangan solusi teknologi untuk mendukung kegiatan operasional. PT Nutech Integrasi dipilih sebagai tempat kerja praktik karena bidang usahanya relevan dengan penerapan teknologi informasi, analisis kebutuhan, dan perancangan solusi sistem.

Dalam pelaksanaan kerja praktik, penulis ditempatkan pada fungsi \textit{Product Specialist} yang berperan dalam eksplorasi, analisis, dan perancangan solusi berdasarkan kebutuhan pelanggan. Kegiatan yang dilakukan meliputi validasi katalog produk, identifikasi permasalahan, eksplorasi teknologi, analisis proses bisnis, serta penyusunan materi \textit{product solution}. Salah satu solusi yang dikaji adalah \textit{Smart Depot Management System}, khususnya teknologi \textit{Pantograph Inspection} dan \textit{Wheel Inspection}. Selain itu, penulis juga melakukan riset \textit{pre-sales} pada sektor energi terbarukan, khususnya Penerangan Jalan Umum Tenaga Surya (PJUTS), untuk mengidentifikasi potensi pasar dan karakteristik calon pelanggan.

Kegiatan kerja praktik tersebut berkaitan dengan fungsi \textit{Product Solution} dalam mendukung pengembangan solusi teknologi dan peluang bisnis perusahaan. Proses ini mencakup analisis kebutuhan operasional, pemodelan alur proses, perancangan dashboard, serta eksplorasi kelayakan solusi dari aspek teknologi dan bisnis. Melalui kegiatan ini, penulis memperoleh pemahaman mengenai proses analisis dan perancangan solusi yang disesuaikan dengan kebutuhan pelanggan serta mendukung arah pengembangan bisnis PT Nutech Integrasi.

\section{Tempat dan Waktu Pelaksanaan}
\label{sec:tempat_waktu}
Kerja praktik dilaksanakan di PT Nutech Integrasi yang beralamat di Jl. Buncit Raya Kav. 99, RT 001/RW 007, Pejaten Barat, Kecamatan Pasar Minggu, Kota Jakarta Selatan, Daerah Khusus Ibukota Jakarta 12510. Pelaksanaan kerja praktik berlangsung selama dua bulan, mulai tanggal 1 Juli 2026 hingga 31 Agustus 2026. Kegiatan kerja praktik dilaksanakan pada hari Senin hingga Jumat, pukul 09.00--18.00 WIB, dengan sistem kerja yang sebagian besar dilakukan secara \textit{Work From Office} (WFO) dan beberapa kegiatan dilaksanakan secara \textit{Work From Anywhere} (WFA).

\section{Ruang Lingkup Kerja Praktik}
Ruang lingkup kerja praktik di PT Nutech Integrasi berfokus pada kegiatan \textit{Product Specialist} dalam mendukung analisis, eksplorasi, dan perancangan konsep solusi teknologi informasi, khususnya pada sektor perkeretaapian serta bidang bisnis pendukung lainnya. Kegiatan ini mencakup validasi produk, analisis proses bisnis, eksplorasi teknologi, penyusunan konsep solusi, serta pengembangan dokumen dan materi pendukung \textit{Product Solution}. Adapun ruang lingkup pekerjaan yang dilakukan selama kerja praktik meliputi:

\begin{enumerate}
    \item Melakukan validasi katalog produk berdasarkan empat pilar bisnis perusahaan, termasuk pemeriksaan aspek legalitas seperti Tingkat Komponen Dalam Negeri (TKDN), Klasifikasi Baku Lapangan Usaha Indonesia (KBLI), dan status hak cipta, serta pemetaan arsitektur teknis solusi \textit{Automatic Fare Collection} (AFC) dari Level 0 hingga Level 4.
    \item Melakukan eksplorasi konsep \textit{Smart Depot Management System}, khususnya teknologi \textit{Wheel Inspection} dan \textit{Pantograph Inspection}, melalui identifikasi permasalahan operasional depo, pemetaan alur fungsional sistem saat ini (\textit{Existing/Confirmed}) dan konsep usulan (\textit{Concept/Proposed}), serta perancangan konsep \textit{Operational Dashboard}.
    \item Menyusun dokumen pendukung \textit{Product Solution}, meliputi draf \textit{product solution deck}, dokumen kebutuhan sistem (\textit{System Requirements}), diagram alur, dan materi pendukung komersial.
    \item Melakukan riset teknologi dan pengumpulan informasi melalui jurnal ilmiah, dokumentasi teknologi, referensi solusi, serta observasi pameran untuk mendukung kegiatan analisis dan pengembangan konsep solusi.
    \item Melakukan eksplorasi awal kelayakan bisnis solusi Penerangan Jalan Umum Tenaga Surya (PJUTS) serta mendokumentasikan hasil diskusi teknis dan komersial bersama mitra PT Jarwinn Solar Panel dalam bentuk \textit{Minutes of Meeting} (MoM).
\end{enumerate}

Pelaksanaan kegiatan kerja praktik memiliki batasan pada tahap analisis, eksplorasi, dan perancangan konsep solusi. Penulis tidak terlibat secara langsung dalam proses implementasi sistem maupun pengembangan perangkat lunak. Selain itu, analisis dilakukan berdasarkan informasi dan data yang dapat diakses selama pelaksanaan kerja praktik. Hasil pekerjaan berupa dokumen analisis, kajian kelayakan awal, rancangan konsep, spesifikasi kebutuhan sistem, diagram, dan materi presentasi yang digunakan sebagai bahan pendukung pengembangan \textit{Product Solution}.

Ruang lingkup tersebut menjadi acuan dalam pelaksanaan kerja praktik serta penyusunan laporan. Penjelasan lebih lanjut mengenai kegiatan, metode analisis, eksplorasi teknologi, dan hasil pekerjaan yang dilakukan selama kerja praktik disajikan pada Bab III.

\section{Sistematika Penulisan}
Sistematika penulisan laporan kerja praktik ini disusun untuk memberikan gambaran mengenai struktur dan alur pembahasan laporan. Laporan kerja praktik terdiri atas enam bab dengan uraian sebagai berikut.

\begin{enumerate}
    \item \textbf{Bab I Pendahuluan}\\
    Bab ini membahas latar belakang pelaksanaan kerja praktik, tempat dan waktu pelaksanaan, ruang lingkup pekerjaan, serta sistematika penulisan laporan.

    \item \textbf{Bab II Organisasi dan Lingkungan Kerja Praktik}\\
    Bab ini menjelaskan profil PT Nutech Integrasi, visi dan misi perusahaan, struktur organisasi, lingkungan kerja, serta jadwal pelaksanaan kerja praktik.

    \item \textbf{Bab III Pelaksanaan Pekerjaan Bulan Juli}\\
    Bab ini menguraikan kegiatan kerja praktik pada bulan pertama, meliputi pengenalan lingkungan kerja, validasi katalog produk, pemetaan solusi \textit{Automatic Fare Collection} (AFC), serta eksplorasi awal teknologi dan solusi perusahaan.

    \item \textbf{Bab IV Pelaksanaan Pekerjaan Bulan Agustus}\\
    Bab ini membahas kegiatan kerja praktik pada bulan kedua, meliputi analisis dan perancangan konsep \textit{Smart Depot Management System}, penyusunan dokumen pendukung \textit{Product Solution}, serta eksplorasi awal kelayakan bisnis solusi Penerangan Jalan Umum Tenaga Surya (PJUTS).

    \item \textbf{Bab V Kegiatan Lainnya}\\
    Bab ini menguraikan kegiatan pendukung selama kerja praktik, meliputi kegiatan internal perusahaan, diskusi dan presentasi, serta aktivitas lainnya di lingkungan kerja.

    \item \textbf{Bab VI Rangkuman}\\
    Bab ini menyajikan rangkuman hasil pelaksanaan kerja praktik, evaluasi diri penulis, serta pengalaman dan pembelajaran yang diperoleh selama kegiatan berlangsung.
\end{enumerate}
